\documentclass[10pt]{article}
\usepackage[a4paper,margin=22mm]{geometry}
\usepackage{amsmath,amssymb,amsthm,mathtools,booktabs,array,enumitem}
\usepackage[T1]{fontenc}
\usepackage{lmodern}
\newtheorem{theorem}{Theorem}
\newtheorem{lemma}{Lemma}
\newtheorem{corollary}{Corollary}
\newcommand{\FINAL}{\mathsf{FINALIZED}}
\newcommand{\ABORT}{\mathsf{ABORTED}}
\newcommand{\REVISED}{\mathsf{REVISED}}
\newcommand{\UPHELD}{\mathsf{UPHELD}}
\newcommand{\DISMISSED}{\mathsf{DISMISSED}}
\newcommand{\MOOT}{\mathsf{MOOT}}
\title{ContestFL Supplementary Material: Extended Proofs and Formal Model}
\author{}
\date{}
\begin{document}
\maketitle

\section{Scope and notation}
For one federated-learning round, let $G_r=(V_r,E_r)$ be the finite acyclic
version graph.  A version commits to its decision key, round identifier, policy,
issuer, result, and the identifiers of every parent version it consumes.  The
ledger state selects at most one current version per active key.  A version is
\emph{terminal} if it is certified or if its renewed challenge window has
closed without a successful challenge.  A checkpoint can be finalized only if
all active keys have current terminal versions, all parent references are
current, the set-containment constraints hold, and the aggregate version binds
the current aggregate-input root and checkpoint.

We separate three assumptions.  Ledger agreement and deterministic contract
execution are denoted by T1.  Binding commitments and authenticated identities
are denoted by T2.  Each registered public verifier or resolver policy has
soundness error $\varepsilon_p$, denoted collectively by T3.  Coverage C1 means
that every policy-invalid provisional version is challenged by an eligible
honest party before its window closes.  Evidence availability C2 means that all
evidence required for successful correction is available before its deadline.

\section{Authenticated-state representation}
% Drop-in specification paragraph for the ContestFL evaluation/design section.
\paragraph{Authenticated-state submission.}
The root-only implementation replaces per-client ledger mappings by four
round-level commitments: a submission root $R_{\mathrm{sub}}$, a decision root
$R_{\mathrm{dec}}$, an aggregate-input root $R_{\mathrm{agg}}$, and the claimed
checkpoint hash $h(w_{r+1})$.  For client identifier $i$, the benchmark derives
a 256-bit key $k_i=H(\mathtt{ContestFL:client}\parallel i)$ and assigns it to the
first $d$ key bits.  A non-empty leaf is
$H(\mathtt{ContestFL:leaf}\parallel k_i\parallel H(v_i))$, and an internal node
is $H(\mathtt{ContestFL:node}\parallel L\parallel R)$; empty subtrees use
height-specific domain-separated default hashes.  The reported experiment uses
$d=32$ and rejects any two distinct keys that map to the same truncated path,
so no collision can silently overwrite a leaf.  This depth is an evaluated
benchmark configuration rather than a collision-free production choice; a
production deployment should use the full key depth or a collision-resolving
authenticated dictionary.

A client-decision challenge reveals $(k_i,v_i)$ and the $d$ sibling hashes.  The
contract reconstructs the path and compares the resulting root with
$R_{\mathrm{dec}}$ inside the same state-changing transaction that records the
challenge.  Consequently, the receipt's gas includes ABI decoding, every
SHA-256 path step, duplicate checking, challenge storage, and event emission.
An aggregate-checkpoint challenge needs no per-client witness because both
$R_{\mathrm{agg}}$ and $h(w_{r+1})$ are already ledger-bound; its evidence hash
commits to the deterministic replay record.  A successful client-decision
revision may install successor decision and aggregate-input roots together with
a successor checkpoint.  A successful aggregate-only revision is constrained
to preserve both roots and may replace only the checkpoint.  Every uncertified
successor reopens a bounded challenge window before finalization.

This implementation is a storage-representation ablation, not a complete
replacement for every materialized ContestFL transition.  In particular,
receipt completeness, omission/injection disputes, challenge flooding,
mooting, and fallback remain evaluated in the materialized prototype unless
explicitly stated otherwise.  Root-only operation also shifts authenticated
state and witness availability off chain: a challenger must obtain the relevant
leaf and path before the challenge deadline.


\section{Extended proofs}

\begin{theorem}[Policy-relative revision soundness]
For an admissible challenge against a policy-valid version, for which all
issuer-obligated evidence is available before the deadline, the probability
that the state machine accepts outcome $\REVISED$ is at most
\[
  \varepsilon_{\mathrm{bind}}+\varepsilon_{\mathrm{auth}}
  +\sum_{p\in\mathcal P_C}\varepsilon_p,
\]
where $\mathcal P_C$ is the set of evidence policies invoked on the resolution
path.
\end{theorem}
\begin{proof}
Let $d$ be the challenged version.  Its identifier binds the round, decision
key, policy identifier, issuer, result, and parent versions.  Challenge
admission binds the reason and counter-evidence to $d$ and authenticates the
challenger.  The state machine can produce $\REVISED$ only if a registered
public verifier accepts the evidence or an authorized resolver certificate
satisfies the registered threshold policy.  Missing challenger evidence maps
to $\DISMISSED$; evidence unavailability attributable neither to the target nor
to the challenger maps to fallback or abort, not to revision.

Assume $d$ is policy-valid and nevertheless revised.  If all bindings and
identities are correct, the accepted verifier/certificate statement is false.
Thus at least one of the following bad events occurs: a commitment collision or
second preimage; an identity, signature, or authorization forgery; an unsound
public verification result; or acceptance of a false resolver certificate.
T2 bounds the first two events and T3 bounds each invoked policy event.  Applying
the union bound proves the claim.  No independence assumption is needed.
\end{proof}

\begin{lemma}[Supersession completeness]
If $G_r$ is acyclic and every version commits to all parent identifiers it
consumes, revising version $d$ and traversing outgoing dependency edges marks
every current descendant of $d$ as superseded.  No marked descendant can remain
in the canonical lineage.
\end{lemma}
\begin{proof}
Order descendants by shortest-path distance from $d$.  Distance-zero contains
$d$, which is marked by the revision transition.  Assume every descendant at
distance at most $q$ is marked.  A descendant $x$ at distance $q+1$ has a parent
$y$ at distance $q$ whose identifier is committed in $x$.  When $y$ is marked,
the traversal visits every outgoing edge of $y$ and marks $x$.  Induction marks
all descendants.  A marked version is not current.  Moreover, any version that
still references a marked parent fails the current-parent predicate even if it
was not yet selected as the current value of its own key.  A replacement can
re-enter the lineage only by committing to the new current parents, so it is a
new version rather than the superseded descendant.
\end{proof}

\begin{theorem}[Lineage-consistent finalization]
Every finalized checkpoint depends only on current terminal versions.  It
cannot depend on a revised target, a descendant superseded by that revision, or
an uncertified replacement whose renewed challenge window remains open.
Under T2--T3 and C1, every version in the finalized lineage is policy-valid
except with the union of the invoked binding and policy failure probabilities.
\end{theorem}
\begin{proof}
Consider the state immediately before finalization.  A revised target is marked
superseded and, by supersession completeness, so are all descendants bound to
its obsolete identifier.  Therefore none satisfies the current-version
predicate.  A replacement without a validating public or quorum certificate is
provisional until its new challenge window closes, so it fails terminality
while that timer is open.  Finalization additionally checks that every active
key has exactly one current terminal version, that every parent reference names
a current version, that the admitted and included sets satisfy the protocol's
nesting constraints, and that the aggregate version binds the current
aggregate-input root and checkpoint.  Hence a finalized aggregate cannot have a
revised, superseded, missing, or provisional ancestor.

For policy validity, suppose a policy-invalid provisional version $x$ occurs in
a finalized lineage.  By C1, an eligible honest party submits an admissible
challenge before $x$ becomes terminal.  Outside the T2--T3 failure events, the
registered resolution either revises $x$ or safely aborts the round.  Revision
removes $x$ and its descendants from the current lineage; abort prevents
finalization.  Both contradict the supposition.  A union bound over the finite
set of bindings and policies invoked in the finalized lineage gives the stated
failure probability.
\end{proof}

\begin{corollary}[Canonical checkpoint uniqueness]
Under T1, two honest federation members cannot accept different finalized
checkpoints for the same round.
\end{corollary}
\begin{proof}
T1 gives every honest member the same totally ordered contract execution.  The
finalization transition is deterministic and can execute at most once for the
round.  Its precondition selects one current terminal aggregate version, whose
committed checkpoint is therefore unique.  A second, distinct finalization
would require either a different ordered state or nondeterministic execution,
contradicting T1.
\end{proof}

\begin{lemma}[Bounded number of challenge windows]
Let levels $1,\ldots,K$ follow a topological layering of $G_r$, let $b_\ell$ be
the non-renewable republication budget shared by all keys at level $\ell$, and
let $F$ be the number of reject-to-admit revisions that create a previously
absent inclusion key.  Then
\[
  E\le 1+\sum_{\ell=1}^{K}b_\ell+F.
\]
\end{lemma}
\begin{proof}
Charge every window after the primary window to exactly one of two finite
resources.  Republishing a key that has already existed consumes one token from
the shared budget of its dependency level; no token can be consumed twice, so
there are at most $\sum_\ell b_\ell$ such windows.  A reject-to-admit revision
may create a previously absent inclusion key.  Charge the first publication of
that key to the corresponding revision.  Since the revised admission is
terminal and the submission receipt is fixed, the same missing key cannot be
created for the first time twice, yielding at most $F$ additional windows.
Certified replacements open no window.  Once a level budget is exhausted, only
a certified fallback or abort is permitted.  Adding the primary window proves
the bound.
\end{proof}

\begin{theorem}[Bounded termination after stabilization]
Assume T1 and finite block-height deadlines.  Let every evidence adapter either
return before its deadline or map timeout to dismissal, certified fallback, or
abort.  Then every round reaches $\FINAL$ or $\ABORT$ in finite time after
network stabilization.
\end{theorem}
\begin{proof}
After stabilization, T1 ensures that every continuously submitted honest
transaction is eventually ordered.  Each submission, decision, challenge,
response, republication, renewed-challenge, fallback, and finalization phase has
a finite block-height deadline.  Therefore every phase either executes its
expected transition or takes its deterministic timeout transition in finite
time.  The preceding lemma bounds the number of challenge windows by the finite
quantity $E_{\max}$.  No unsupported challenge consumes a replacement token,
and exhaustion cannot reopen a provisional window.  Consequently the recovery
loop cannot execute indefinitely.  If a finalizable state is reached, the
finalization transition terminates the round; otherwise a timeout, unavailable
fallback, or budget exhaustion takes the round to $\ABORT$.
\end{proof}

\begin{theorem}[Bounded unsupported-challenge workload]
Suppose each adversarial identity may open at most $\theta$ challenges per
window, duplicate target/reason pairs are merged, there are $f$ adversarial
identities, $n_e$ provisional decisions and $|\mathcal P|$ active reason codes
in epoch $e$, and at most $p$ adapters execute concurrently.  If one unsupported
challenge takes at most $\Delta_{\max}$ to resolve, its additional resolution
delay is bounded by
\[
D_{\mathrm{adv}}^{(e)}\le
\left\lceil\frac{\min(f\theta,n_e|\mathcal P|)}{p}\right\rceil
\Delta_{\max}.
\]
\end{theorem}
\begin{proof}
The identity quota permits at most $f\theta$ submissions.  Duplicate merging
also limits distinct target/reason pairs to $n_e|\mathcal P|$, so at most the
minimum of these quantities reaches an adapter.  Scheduling that many jobs on
$p$ parallel workers requires at most the displayed number of batches, each of
duration at most $\Delta_{\max}$.  Unsupported challenges do not revise state,
consume replacement budgets, or open renewed windows.  Summing over the bounded
number of epochs gives a finite total adversarial delay.
\end{proof}

\section{Finite-state TLA+ model}
The accompanying \texttt{ContestFL.tla} is a deliberately finite abstraction of
the recovery lifecycle.  It models provisional and certified checkpoints,
valid and false challenges, supersession, bounded retries, certified fallback,
abort, and at-most-once finalization.  It does not model cryptography, concrete
Merkle trees, the complete multi-key dependency DAG, wall-clock performance, or
unbounded populations.  Those aspects are covered by the mathematical argument
or measured implementation rather than inferred from bounded model checking.

\begin{table}[h]
\centering
\caption{Mapping between protocol claims and checked TLA+ predicates.}
\begin{tabular}{@{}p{0.42\linewidth}p{0.50\linewidth}@{}}
\toprule
Protocol claim & TLA+ predicate \\
\midrule
Retry budgets cannot be exceeded & \texttt{RetryBound} \\
At most one checkpoint is finalized & \texttt{AtMostOneFinalization} \\
A superseded checkpoint is never finalized & \texttt{NoSupersededFinalization} \\
A provisional replacement is never finalized & \texttt{NoProvisionalFinalization} \\
Finalization selects the current checkpoint & \texttt{CanonicalFinalization} \\
With challenge coverage, finalized checkpoint is valid & \texttt{PolicyValidityUnderCoverage} \\
Terminal states have no open challenge & \texttt{TerminalHasNoOpenChallenge} \\
Every behavior eventually finalizes or aborts & \texttt{Termination} \\
\bottomrule
\end{tabular}
\end{table}

The model is checked in two configurations.  \texttt{MC\_Coverage.cfg} enables
coverage and checks policy-valid finalization.  \texttt{MC\_Structural.cfg}
disables coverage and checks only structural safety and termination, allowing an
invalid but unchallenged checkpoint to finalize.  The artifact runner records
the exact TLC version, complete console output, generated and distinct state
counts, graph depth, elapsed time, and both configuration files.  The generated
\texttt{TLC\_RESULTS.tex} can be inserted below after executing the artifact.

\IfFileExists{TLC_RESULTS.tex}{\input{TLC_RESULTS.tex}}{%
\paragraph{TLC execution placeholder.}
Run \texttt{formal/run\_tlc.sh}; then copy the generated
\texttt{TLC\_RESULTS.tex} next to this document.  No state counts are asserted
in this source before that exhaustive run succeeds.}

\section{Reproducibility boundary}
The proofs are parameterized; TLC checks the finite constants in the supplied
configuration files.  Conversely, the implementation measurements establish
cost and behavior for the evaluated workloads but cannot prove all executions.
The paper should therefore report the two forms of evidence separately: theorem
statements for parameterized safety/termination, bounded exhaustive state counts
for the formal abstraction, and empirical gas/latency/replay measurements for
the implementation.

\end{document}
